\subsection{The common-stabilizer measure at the arithmetic reference lattice}
\label{sec:flag-reference-haar}

We now use the normal form of Section~\ref{sec:flag-adapted-normal-form}
to evaluate the reference covolume in the \emph{local orbital}
measure. This is not the Hilbert--Schmidt measure on the Lie
algebra used in the earlier unipotent-volume calculation.

Write \(\rho=\operatorname{rd}(E)=2v\), \(f=[F:\mathbb Q]\),
\(A_a=\mathfrak d^{a-r}\), and
\[
       L_v=(8v)^{2(r+1)^2f}.
\]
Here \(v\ge1\), and \(N_E\mathfrak d=\rho^{2f}\), as in the
original arithmetic data. Continue with (US1)--(US3). In
particular the middle space \(W+Ez\) has rank \(k+1\), and its
lattice \(J\) in the reference construction is self-dual by
Lemma~\ref{lem:flag-fixed-index-constraint}. No self-duality of the
whole reference lattice is needed.

\begin{theorem}[Reference and actual common-stabilizer covolumes]
\label{thm:flag-reference-haar}
For an actual admitted rational normal form with common tail of
length \(r-j\), let \(R\) be the common flag unipotent group.
Use the original real model of this normal form: the first
isotropic vectors have coordinates \(\rho^{r-a}e_a^{\rm std}\)
and their duals have coordinates
\(\rho^{a-r}f_a^{\rm std}\). Let \(\lambda_R\) be product
Lebesgue Haar measure in the original block-upper, middle and
skew coordinates, i.e. the product of the source's
\(\mathrm{lebR0}\) measures with scalar \(1\).
Then
\[
 \frac{v^{k(r-j)f}}{L_v}
 \ \le\ \operatorname{covol}_{\lambda_R}\Gamma_{\rm ref}
 \ \le\ v^{k(r-j)f} L_v ,
 \tag{RH1}
\]
where \(\Gamma_{\rm ref}=R(E)\cap\operatorname{Aut}(\Lambda_{\rm ref})\).
The actual lattice group
\(\Gamma=R(E)\cap\operatorname{Aut}(\Lambda)\) satisfies
\[
 \operatorname{covol}_{\lambda_R}\Gamma
       \le \omega_N(\Lambda)v^{k(r-j)f} L_v ,
 \qquad
 \omega_N(\Lambda)=
 \frac{[\Gamma_{\rm ref}:\Gamma_{\rm ref}\cap\Gamma]}
      {[\Gamma:\Gamma_{\rm ref}\cap\Gamma]}.
 \tag{RH2}
\]
In particular the covolumes are finite and strictly positive.
\end{theorem}

\begin{proof}
Use the orthogonal splitting (US1), and put \(I=X\cap X'\),
\(I^*=\langle f_a:a\ge j\rangle\), \(M_E=W+Ez\).
We write matrices with the conjugate-linear first argument
convention; this merely conjugates the entries of the matrices
in (US1). On \(X_j+X'_j\), an element of \(R\) acts on
\(X_j\) by an upper unitriangular matrix \(U\).
Its action on \(X'_j\) is forced by the cross-pairing matrix
\(P\) to be \(P^{-1}(U^*)^{-1}P\), where \(*\) is conjugate
transpose. Both matrices must be upper unitriangular.
Since \(P\) is a direct sum of signed antidiagonals, reversing
indices within each block takes the inverse-transpose upper
condition into the same condition on the first block. Between
distinct blocks the inverse-transpose is lower triangular in
block order; hence its off-block entries must vanish.
Thus \(U\) is precisely block diagonal upper unitriangular
on the leading blocks.

On the orthogonal complement \(I+M_E+I^*\), the common
stabilizer is the usual full flag unipotent group for \(I\).
It acts trivially on \(I^\perp/I\) in this complement.
There is no further constraint from the second flag because its
tail is the first tail multiplied by a fixed norm-one scalar.
In particular, the middle action is the identity, and not
the compact stabilizer of \(z\). If \(H_M\) denotes the matrix of
the middle form, the radical of this parabolic has matrices
\[
 p(B,S)=
 \begin{pmatrix}
 1&-B^*H_M&S-\tfrac12 B^*H_MB\\
 0&1&B\\0&0&1
 \end{pmatrix},
 \qquad S^*=-S,
\]
and the remaining upper unitriangular tail acts on \(I\) and
contragrediently on \(I^*\). This follows by multiplication
against the Hermitian matrix of the orthogonal splitting.
Conversely all these matrices preserve both flag filtrations
and act trivially on both isotropic quotients, so the description
is exact. Reversal of the dual ordering gives the coordinates
\(\mathrm{rEl}\) in the source. Let \(n_1\) be the number of upper
entries in the leading blocks and the tail, and \(n_3\) the number
of strictly upper entries in the skew tail. Set \(s=r-j\).
Then \(n_1,n_3\le r^2\) and \(s\le r\).

In the reference lattice, each first isotropic line has
coefficient ideal \(A_a\), each second leading line has the
same ideal, and the common dual tail has ideals \(A_a^{-1}\).
The middle lattice is \(J\). Thus every upper entry with
indices \(a<b\) is in \(A_a A_b^{-1}=\mathfrak d^{a-b}\).
These conditions suffice: they are exactly integrality on the
first upper module, including its inverse, and the paired
leading action reverses the entry indices inside the same
block. The two indices are replaced by their reversed pair
with the same difference \(a-b\); the block scalar cancels,
and the column signs are units. On the tail the dual action
is integral by perfect duality. Hence pure upper matrices
lift every allowed upper coordinate.

For the identity-upper subgroup, the images of
\(A_b^{-1}f_b\), \(b\ge j\), force \(B_b\in A_b J\).
Conversely every such family lifts integrally.
There exists \(w_0\in\mathcal O_E\) with
\(w_0+\bar w_0=1\): the trace ideal is the unit ideal locally
at every finite prime by unramifiedness (also at split primes
and above two), and therefore globally. Replace the upper-right
block \(-\tfrac12 B^*H_MB\) by \(-w_0 B^*H_MB\). Their difference
is skew-Hermitian. The entries of \(B^*H_MB\) are in
\(A_a A_b\) because \(J\) is self-dual and the ideals are
conjugation invariant. The other off-diagonal block
\(-B^*H_M\) maps \(J\) into the required first tail module.
These conditions give an integral unipotent automorphism,
including its inverse. The central kernel \(B=0\) consists exactly
of skew matrices with \(S_{ab}\in A_a A_b\). In particular its
diagonal consists of anti-invariant elements of \(A_a^2\).
This proves both necessity and surjectivity of the projected
integral groups; it does not assert that arbitrary integral
\(B,S\) give an integral matrix with the factor \(1/2\).

The group law on the successive upper-entry quotients (ordered
by distance from the diagonal), the middle columns and the
skew kernel is additive, with triangular translation in each
fibre. The upper matrices act on the middle and skew spaces
with determinant one. Thus product coordinate Lebesgue
measure is Haar and covolume is the product of the successive
additive-lattice covolumes. For a direct justification, reduce
each quotient into a half-open fundamental parallelepiped
and lift the reducing translation by the just-proved
surjectivity. The remaining coordinate changes are
translations or triangular determinant-one transformations.
This gives a Borel fundamental set, with its volume computed
by iterated integration; no rectangular integrality claim
for the nonabelian group is used.

The complex-embedding covolume of a fractional ideal \(A\) is
\(v^f N_E(A)\). In the prescribed real coordinates an upper
entry is multiplied at each place by \(\rho^{b-a}\), and
an off-diagonal skew entry by \(\rho^{2r-a-b}\).
The corresponding ideal norms cancel the real Jacobians
\(\rho^{2f(b-a)}\) and \(\rho^{2f(2r-a-b)}\).
Each of these coordinates consequently has covolume \(v^f\).

The middle column at index \(b\) is multiplied by
\(\rho^{r-b}\). For an integral self-dual Hermitian lattice of
rank \(d\), its embedding covolume in the standard Hermitian
space of determinant modulus one is \(v^{fd}\).
Indeed take a pseudo-basis over the Dedekind domain and compute
the determinant of its real embedding; the absolute norm of
the Hermitian determinant times the ideal factors is one by
self-duality. Multiplication by \(A_b\) changes the volume by
\(N_E(A_b)^d\). The rescaling of the \(d=k+1\) complex
coordinates cancels this factor exactly. Each middle column
therefore has covolume \(v^{f(k+1)}\), even if \(J\) is not a
free \(\mathcal O_E\)-module.

For a skew diagonal coordinate at \(a\), take the full complex
lattice
\[
 \mathcal L=\rho^{2(r-a)}
            \bigl(\sigma_\tau(A_a^2)\bigr)_\tau
             \subset \mathbb C^f .
\]
It is invariant under conjugation, has covolume \(v^f\), and
every nonzero vector has sup norm at least one: its product
of squared coordinate moduli is at least one by the fractional
ideal norm inequality. The intersections with the real and
imaginary subspaces are full lattices of rank \(f\), because
their direct sum contains \(2\mathcal L\). Each has covolume
at least one by Minkowski's convex body theorem applied to
the open unit cube. Their product is at most \(4^f v^f\),
since their direct sum has index at most \(2^{2f}\) in
\(\mathcal L\). The imaginary covolume therefore lies between
\(1\) and \((4v)^f\), in particular between \((8v)^{-f}\)
and \((8v)^f\).

The product formula for the reference group now gives
\[
 v^{f(n_1+n_3+(k+1)s)}
       \prod_{a=j}^{r-1} V_a ,
 \qquad (8v)^{-f}\le V_a\le(8v)^f.
\]
Because \(v\ge1\), \(n_1+n_3+2s\le 2r^2+2r\),
this implies (RH1), also for an empty tail. The two
integral stabilizers are commensurable: choose an integer
\(d>0\) with \(d\Lambda\subseteq\Lambda_{\rm ref}\subseteq
d^{-1}\Lambda\), then the kernel of the finite action on
\(d^{-1}\Lambda/d\Lambda\) preserves both lattices; interchange
the lattices for the opposite index. The finite-index formula
for the common Haar covolume proves (RH2).
\end{proof}

\paragraph{Real realization and the two orientations.}
The real coordinates in this proof come from the actual
rational frame, not from a new local flag orbit. Complete its
first flag to a hyperbolic basis by taking, for the first \(j\)
duals, the dual basis inside the paired leading space \(X'_j\),
and for the remaining duals the given \(f_a\). They are all
isotropic, and perpendicular to the last vector. Complete the
second flag in the same way, using \(c f_a\) on the tail.
A rational isometry sends the first adapted hyperbolic basis
to the original hyperbolic basis and sends its last vector
to the original last vector. On the middle orthogonal
complement take the canonical isometry through \(X^\perp/X\).
It sends the last vector correctly because adaptation preserves
its class there. The resulting isometry acts trivially on all
graded flag spaces and on \(X^\perp/X\), hence belongs to the
original flag unipotent group. Apply the same construction
to the second adapted basis pulled back by the original
return \(\gamma^{-1}\). Call the two isometries \(m_2,m_A\).
In the original complex realization \(\phi_\tau\) this gives
\[
 \Phi_\tau=\phi_\tau m_2,\qquad
 g_0=m_A\gamma^{-1}m_2^{-1},\qquad
 \gamma^{-1}=m_A^{-1}g_0m_2
\]
(with the rational maps represented by their complex matrices
in the latter two equations). The leading block moduli of
\(g_0\) are constant on each antidiagonal: the original
pairing is multiplied by
\(\rho^{-(r-a)-(r-b)}\), constant when \(a+b\) is fixed.
On the common tail and its duals \(g_0\) acts by \(c^{-1}\)
at this place, and it maps the last vector by this same scalar.
Thus these are exactly the source's \(\mathrm{IsFrameOf}\),
\(\mathrm{IsNormal0}\) and \(\mathrm{lebR0}\) coordinates.
The measure on the original stabilizer is pushed back by
\(m_2^{-1}\). For the inverse orientation use the conjugation
and equal-scalar calculation in the proof of
Theorem~\ref{thm:actual-local-degenerate-bound}; it is measure
preserving in the stipulated normal-form measures. Consequently
(RH2) is the actual \(\mathrm{CovTC}\) bound. It uses the
relative weight of this adapted form, before the later
infimum over adapted forms.

For the reference word's base lattice take \(j=0\), identical
flags, and the given base hyperbolic splitting. Then its
reference lattice is the base lattice itself and the above
Lebesgue measure is \(\mathrm{lebN}\). The original membership
calculation identifies the lattice group with the unipotent
period lattice (Theorem~\ref{global-gram-haar-covolume}).
In the normalized local coordinates at \(\theta=1\), the
original unipotent measure is \(\kappa\,\mathrm{lebN}\), where
\[
 \kappa=(2v)^{-rk}2^{r(k+1)+r(r-1)}
        =2^{r^2}v^{-rk}.
\]
Hence the product-coordinate unipotent covolume \(V_N\)
in the original period normalization satisfies
\[
             0<V_N<\infty,\qquad V_N L_v\ge1 .
 \tag{RH3}
\]
Indeed multiply the lower bound (RH1) at \(j=0\) by
\(\kappa^f\). This is the source's \(\mathrm{CovNC}\).
In particular all measure factors used in the subsequent
amplitude estimate have the required sign and finiteness.
